\section{The exact-WKB thesis and the computational spine}
\label{sec:computational-spine}
% BEGIN CURATED SUBJECT INDEX
\index{Wronskian}
\index{connection coefficient}
\index{boundary condition!incoming and outgoing}
\index{resonance!residue}
\index{Borel summation}
\index{Stokes automorphism}
\index{Riccati equation}
\index{exact WKB}
\index{monodromy}
\index{exceptional point}
% END CURATED SUBJECT INDEX

The organization of this book follows a claim that is stronger than the
statement that exact WKB is useful for finding resonance energies.  The claim
is that a resonance problem is most economically organized around its
\textit{global analytic connection datum}.  The wave function is first known
only through local asymptotic representations.  Exact WKB tells us how to
construct these representations and how they jump; scattering theory tells
us which global coefficient is physically the incoming amplitude; complex
scaling and rigged spaces tell us how the resulting pole functional may be
represented.  This order prevents a useful representation from being
mistaken for the definition of the resonance.

For a second-order equation
\begin{equation}
  \left[-\hbar^2\frac{\rmd^2}{\rmd x^2}+Q(x,E)\right]\psi(x,E)=0,
  \label{eq:spine-equation}
\end{equation}
the calculation begins with the spectral curve
\begin{equation}
  y^2=Q(x,E).
  \label{eq:spine-spectral-curve}
\end{equation}
The two sheets of $y=\sqrt{Q}$ label the leading WKB branches.  The Riccati
recursion promotes them to formal all-order solutions, Borel summation gives
sectorial analytic functions, and Stokes automorphisms relate neighboring
sectors.  After transporting a basis from one asymptotic region to another,
one obtains a connection matrix
\begin{equation}
  \bm{\vsub{a}{R}}(E)
  =\mathcal M(E)\bm{\vsub{a}{L}}(E),
  \qquad
  \mathcal M(E)=\mathcal S_N(E)\cdots\mathcal S_2(E)\mathcal S_1(E),
  \label{eq:spine-global-connection}
\end{equation}
where each $\mathcal S_j$ is a local Stokes, branch-cut, or propagation
matrix.  Wronskian preservation constrains $\det\mathcal M$ and supplies an
immediate error check.

The boundary condition selects a scalar minor or matrix block of
$\mathcal M$.  We denote the coefficient multiplying the prohibited incoming
component by $\Cincoming(E)$.  The resonance condition is
\begin{equation}
  \Cincoming(\resonantE)=0,
  \qquad
  \resonantE=E_*-\frac{\rmi}{2}\Gamma .
  \label{eq:spine-resonance-condition}
\end{equation}
For a simple zero, a scattering amplitude has the local form
\begin{equation}
  S(E)=\frac{\mathcal N(E)}{\Cincoming(E)},
  \qquad
  \Res_{E=\resonantE}S(E)
  =\frac{\mathcal N(\resonantE)}
  {\partial_E\Cincoming(\resonantE)}.
  \label{eq:spine-residue}
\end{equation}
Thus the pole, decay width, and residue are successive pieces of information
in one calculation.  A multiple zero anticipates an exceptional point.  A
branch point in $E$ anticipates a threshold.  A singular endpoint changes
the monodromy entering $\mathcal M$.  These are not separate subjects added
to resonance theory; they are different failures of the simplest connection
problem and therefore have a common diagnostic language.

\subsection{What is novel in this ordering}
\label{subsec:novel-ordering}
% BEGIN CURATED SUBJECT INDEX
\index{boundary condition!incoming and outgoing}
\index{resonance!residue}
\index{Gamow--Siegert state}
\index{exact WKB}
\index{monodromy}
\index{Zel'dovich regularization}
\index{continuum level density}
\index{exceptional point}
\index{holonomy}
\index{quasinormal mode}
% END CURATED SUBJECT INDEX

The exact-WKB viewpoint produces five conceptual shifts that guide the rest
of the book.

\begin{enumerate}
  \item \textbf{The pole condition precedes the pole state.}
  A Gamow--Siegert function grows on the real axis, but the vanishing incoming
  coefficient is already well defined by analytic continuation.  One should
  therefore establish Eq.~\eqref{eq:spine-resonance-condition} before asking
  in which space the associated state lives.

  \item \textbf{Regularizations are representations of one analytic
  continuation.}
  Zel'dovich damping, contour rotation, and a rigged-space pairing appear
  different when applied directly to a divergent wave function.  Exact WKB
  shows what they must preserve: the same continued branch, connection
  coefficient, and residue.  Their equivalence is consequently tested by
  analytic deformation invariance rather than by a formal equality of
  divergent integrals.

  \item \textbf{The continuum and the resonance belong to the same global
  problem.}
  The continuum is not a disposable background left after a pole has been
  isolated.  It supplies the boundary values, cuts, Stokes sectors, and
  completeness contour from which the pole is continued.  This is why
  continuum level density and pole residue are natural observables alongside
  the pole position.

  \item \textbf{Monodromy supplies a renormalization condition.}
  At a radial singularity or a continuum-normalized exceptional point, an
  auxiliary regulator enters local normalization.  Requiring the global
  monodromy or holonomy to be independent of that auxiliary choice produces
  a renormalization-group-type condition.  The invariant is global even when
  the counterterm is local.

  \item \textbf{Black-hole QNMs are not an analogy but another connection
  calculation.}
  Ingoing and outgoing horizon conditions select the prohibited coefficient
  exactly as in Eq.~\eqref{eq:spine-resonance-condition}.  The difference lies
  in the singularities and asymptotic regions of the spectral curve, not in a
  new definition of resonance.
\end{enumerate}

\subsection{The calculation cycle used in every laboratory}
\label{subsec:calculation-cycle}
% BEGIN CURATED SUBJECT INDEX
\index{topology}
\index{spectral measure}
\index{resolvent}
\index{Wronskian}
\index{connection coefficient}
\index{boundary condition!incoming and outgoing}
\index{branch cut}
\index{resonance!residue}
\index{scattering!Coulomb}
\index{infrared problem}
% END CURATED SUBJECT INDEX

Each worked problem follows the same cycle.  The reader should use it as an
algorithm rather than as a list of topics.

\begin{enumerate}
  \item Specify the differential operator, its domain before continuation,
  the time convention, and every asymptotic channel.

  \item Draw or determine the spectral curve, turning points, singular
  points, branch cuts, and relevant Stokes graph.

  \item Generate the Riccati coefficients and decide which formal quantities
  require Borel summation or another resummation.

  \item Choose normalized local bases and multiply the connection matrices in
  path order.  Check the Wronskian or determinant after every multiplication.

  \item Identify $\Cincoming(E)$ from the physical boundary
  condition.  Locate its zeros on a declared sheet and Borel side.

  \item Differentiate the connection datum to obtain residues or
  normalization data, and retain the non-pole part when a response function
  is required.

  \item Only now choose a convenient realization: a complex-scaled $L^2$
  eigenproblem, a Zel'dovich pairing, a rigged-space functional, a momentum
  bin, or a direct-integral spectral density.

  \item Vary the representation data---scaling angle, contour, basis size,
  bin width, regulator, or resummation order---while holding the analytic
  invariant fixed.  Stability is part of the identification of a resonance,
  not a cosmetic numerical test.
\end{enumerate}

\begin{longtable}{@{}p{0.12\linewidth}p{0.33\linewidth}p{0.47\linewidth}@{}}
\toprule
stage & new input & calculational capability \\
\midrule
\endfirsthead
\toprule
stage & new input & calculational capability \\
\midrule
\endhead
I-A & topology, domains, and spectral measures & distinguish norm, weak,
strong, and resolvent limits and identify the continuum of a self-adjoint
Hamiltonian \\
\addlinespace
I-B & multiplicity fibers, operator algebras, and compact kernels & pass
between spectral direct integrals, commutants, Fredholm determinants, and
controlled finite model spaces \\
\addlinespace
II & asymptotic channels, flux, and Jost data & construct the physical
scattering operator and identify when Coulomb, many-charge, QED, or gauge
sectors require a different asymptotic representation \\
\addlinespace
III & turning points and Stokes data & construct $\Cincoming$, find a
resonance, and calculate its width and residue \\
\addlinespace
IV & analytic deformation and test-space topology & represent the same pole
by continued resolvents, complex scaling, regularized norms, rigged-space
functionals, and continuum response \\
\addlinespace
V & a regular singular endpoint & include Frobenius indices, the Langer term,
and radial monodromy in exact quantization \\
\addlinespace
VI & few-body channels and model-space projections & derive CDCC and connect
pseudostate convergence, complex-scaled smoothing, and nuclear observables \\
\addlinespace
VII & a multiple zero and a continuum partner & calculate Jordan data,
self-orthogonality, sheet exchange, and regulator-independent holonomy \\
\addlinespace
VIII & horizons and tortoise coordinates & derive Regge--Wheeler-type master
equations and calculate black-hole QNMs and continuum response \\
\bottomrule
\end{longtable}

The short-range one-channel problem is the training ground, not the presumed
universal asymptotic form.  Coulomb and QED examples later show that the first
step of the cycle can itself fail: if the comparison dynamics or asymptotic
representation is wrong, no refinement of the later pole calculation repairs
it.  This is why long-range scattering and operator-algebraic sectors belong
inside the computational spine of the theory.
